\begin{thebibliography}{Cha24b}
\bibitem[Arn61]{Arnold-61} V. I. Arnol’d – “Some remarks on flows of line elements and frames”, Soviet Math. Dokl. 2 (1961), p. 562–564.
\bibitem[CL]{Cekic-Lefeuvre-24} M. Cekić \& T. Lefeuvre – “Semiclassical analysis on principal bundles”, arXiv:2405.14846.
\bibitem[Cha25]{Charles--23} L. Charles , “Resolvents of Bochner Laplacians in the semiclassical limit”, Comm. Partial Differential Equations 50 (2025), no. 7, p. 899–930.
\bibitem[Cha00]{Charles-00} L. Charles – “Aspects semi-classiques de la quantification géométrique”, PhD Thesis, Université Paris Dauphine-PSL, 2000.
\bibitem[Com86]{Comtet-86} A. Comtet – “On the Landau levels on the hyperbolic plane”, Ann. Physics 173 (1986), p. 185–209.
\bibitem[DJN22]{Dyatlov-Jin-Nonnenmacher-22} S. Dyatlov, L. Jin \& S. Nonnenmacher – “Control of eigenfunctions on surfaces of variable curvature”, J. Amer. Math. Soc. 35 (2022), no. 2, p. 361–465.
\bibitem[DZ19]{Dyatlov-Zworski-book} S. Dyatlov \& M. Zworski – Mathematical theory of scattering resonances, Grad. Stud. Math., vol. 200, American Mathematical Society, Providence, RI, 2019.
\bibitem[Gin96]{Ginzburg-96} V. L. Ginzburg – “On the existence and non-existence of closed trajectories for some Hamiltonian flows”, Math. Z. 223 (1996), no. 3, p. 397–409.
\bibitem[IL94]{Iengo-Li-94} R. Iengo \& D. Li – “Quantum mechanics and quantum Hall effect on Riemann surfaces”, Nuclear Phys. B 413 (1994), no. 3, p. 735–753.
\bibitem[Cha24a]{Landau1} L. Charles , “Landau levels on a compact manifold”, Ann. H. Lebesgue 7 (2024), p. 69–121.
\bibitem[Cha24b]{Landau2} L. Charles , “On the spectrum of nondegenerate magnetic Laplacians”, Anal. PDE 17 (2024), no. 6, p. 1907–1952.
\bibitem[Lef26]{Lefeuvre-book} T. Lefeuvre – Microlocal analysis in hyperbolic dynamics and geometry, Cours Spécialisés, vol. 32, Société Mathématique de France, Paris, 2026.
\bibitem[Non19]{Nonnenmacher-notes} S. Nonnenmacher – “Introduction to semiclassical analysis”, Lecture notes, available at https://www.imo.universite-paris-saclay.fr/~stephane.nonnenmacher/enseign/ Cours-Semiclassical-Analysis2019-lecture1-8.pdf, 2019.
\bibitem[TP06]{Tejero-06} C. Tejero Prieto – “Holomorphic spectral geometry of magnetic Schrödinger operators on Riemann surfaces”, Differential Geom. Appl. 24 (2006), no. 3, p. 288–310.
\bibitem[Wei77]{Weinstein-77} A. Weinstein – “Asymptotics of eigenvalue clusters for the Laplacian plus a potential”, Duke Math. J. 44 (1977), p. 883–892.
\bibitem[Zwo12]{Zworski-12} M. Zworski – Semiclassical analysis, Grad. Stud. Math., vol. 138, American Mathematical Society, Providence, RI, 2012.
\bibitem[Bur90]{Burger-90} M. Burger – “Horocycle flow on geometrically finite surfaces”, Duke Math. J. 61 (1990), no. 3, p. 779–803.
\bibitem[Fur73]{Furstenberg-73} H. Furstenberg – “The unique ergodicity of the horocycle flow”, in Recent advances in topological dynamics (New Haven, Conn., 1972), Lect. Notes in Math., vol. 318, Springer, Berlin-New York, 1973, p. 95–115.
\bibitem[FF03]{Flaminio-Forni-03} L. Flaminio \& G. Forni – “Invariant distributions and time averages for horocycle flows”, Duke Math. J. 119 (2003), no. 3, p. 465–526.
\bibitem[Pat99]{Paternain-99} G. P. Paternain – Geodesic flows, Progress in Math., vol. 180, Birkhäuser Boston, Inc., Boston, MA, 1999.
\end{thebibliography}
